% Source: https://artofproblemsolving.com/community/c3h2224698_how_do_you_solve_this_type_of_number_theory_problem
% Stable identity: aops-topic-2224698-q01
\textbf{(d) first: the method.} If $n^2+x=m^2$ with $m>n\ge1$, then
\[x=m^2-n^2=(m-n)(m+n).\]
Put $d=m-n$ and $e=m+n$. Then $de=x$, $d<e$, and $d$, $e$ have the same parity (their sum $2m$ is even). Conversely, any such factorisation gives $m=\dfrac{e+d}{2}$ and $n=\dfrac{e-d}{2}$, a positive integer. So the solutions $n$ correspond exactly to factorisations $x=de$ with $d<e$ and $d\equiv e\pmod2$, via $n=\dfrac{e-d}{2}$. In particular, if $x\equiv2\pmod4$ there are no solutions (one factor would be odd and the other even).

\textbf{(a)} $10=1\cdot10=2\cdot5$: both factorisations mix parities ($10\equiv2\pmod4$), so there is \textbf{no} such $n$.

\textbf{(b)} $56=1\cdot56=2\cdot28=4\cdot14=7\cdot8$. The same-parity pairs are $(2,28)$ and $(4,14)$, giving $n=13$ and $n=5$. So $\boxed{n=5\text{ or }13}$: $25+56=81=9^2$ and $169+56=225=15^2$.

\textbf{(c)} $45=1\cdot45=3\cdot15=5\cdot9$, all odd pairs, giving $n=22,6,2$. So $\boxed{n=2,\ 6\text{ or }22}$: $4+45=7^2$, $36+45=9^2$ and $484+45=23^2$.
